% p213_2; source PDF page 213; manually verified editable table.
\begin{tikzpicture}[x=1cm,y=-1cm,inner sep=0pt,outer sep=0pt]
\path[use as bounding box] (0,0) rectangle (16.000,5.050);
\draw[line width=.4pt] (0.000,0.000) rectangle (2.800,0.850);
\node[align=center,text width=2.540cm,font={\linespread{1}\fontsize{10}{13.50}\selectfont\bfseries}] at (1.400,0.425) {种类};
\draw[line width=.4pt] (2.800,0.000) rectangle (8.900,0.850);
\node[align=center,text width=5.840cm,font={\linespread{1}\fontsize{10}{13.50}\selectfont\bfseries}] at (5.850,0.425) {事项};
\draw[line width=.4pt] (8.900,0.000) rectangle (16.000,0.850);
\node[align=center,text width=6.840cm,font={\linespread{1}\fontsize{10}{13.50}\selectfont\bfseries}] at (12.450,0.425) {机关};
\draw[line width=.4pt] (0.000,0.850) rectangle (2.800,2.650);
\node[align=left,text width=2.540cm,font={\linespread{1}\fontsize{10}{13.50}\selectfont}] at (1.400,1.750) {特别减轻处罚};
\draw[line width=.4pt] (2.800,0.850) rectangle (8.900,2.650);
\node[align=left,text width=5.840cm,font={\linespread{1}\fontsize{10}{13.50}\selectfont}] at (5.850,1.750) {第63条规定，犯罪分子没有法定减轻处罚情节，欲减轻处罚的};
\draw[line width=.4pt] (8.900,0.850) rectangle (16.000,2.650);
\node[align=left,text width=6.840cm,font={\linespread{1}\fontsize{10}{13.50}\selectfont}] at (12.450,1.750) {经最高人民法院核准，可以减轻处罚};
\draw[line width=.4pt] (0.000,2.650) rectangle (2.800,5.050);
\node[align=left,text width=2.540cm,font={\linespread{1}\fontsize{10}{13.50}\selectfont}] at (1.400,3.850) {特别假释};
\draw[line width=.4pt] (2.800,2.650) rectangle (8.900,5.050);
\node[align=left,text width=5.840cm,font={\linespread{1}\fontsize{10}{13.50}\selectfont}] at (5.850,3.850) {第81条第1款规定，有期徒刑欲假释，须执行原判刑期1/2以上；无期徒刑欲假释，须实际执行13年以上};
\draw[line width=.4pt] (8.900,2.650) rectangle (16.000,5.050);
\node[align=left,text width=6.840cm,font={\linespread{1}\fontsize{10}{13.50}\selectfont}] at (12.450,3.850) {如有特殊情况，经最高人民法院核准，可以不受前述执行刑期的限制};
\end{tikzpicture}
